%%
%% Part 1 – Title, Abstract, Introduction and Methodology
%%
%% This file constitutes the first chapter of our modular report.  It
%% introduces the motivation behind generative modelling, states the
%% problem and outlines the methodological framework used in the
%% remainder of the report.  Each section retains the core content
%% from the original assignment while adding narrative glue and
%% clarifying remarks.  Feel free to edit or extend these
%% explanations to improve readability without altering the scientific
%% substance.

\section*{Abstract}
\addcontentsline{toc}{section}{Abstract}
\label{sec:abstract}

The synthesis of realistic financial time series remains a
significant challenge due to non‑stationary dynamics, high
noise‑to‑signal ratios and the presence of heavy‑tailed distributions.
Deep generative models offer powerful tools to address this challenge
by learning complex probability distributions from data.  In this
report we investigate the transition from discrete‑time \emph{Denoising
Diffusion Probabilistic Models} (DDPM) to continuous‑time \emph{Flow
Matching} (FM) frameworks.  We demonstrate the mathematical
equivalence between the probability flow ordinary differential
equations (ODEs) of diffusion and the deterministic vector fields of
Flow Matching.  Using SPY exchange traded fund (ETF) log‑returns as a
benchmark, we implement a transformer‑based Flow Matching model that
utilises linear optimal transport paths.  Our results indicate that
while both models effectively capture market volatility clustering,
Flow Matching offers superior sampling efficiency and stability.  We
evaluate the generated data using the \emph{Sliced Wasserstein
Distance} (SWD) and temporal autocorrelation metrics, revealing high
fidelity in reproducing first‑ and second‑order moments, albeit with
noted challenges in capturing extreme leptokurtic tail behaviour.

\section{Introduction}
\label{sec:introduction}

\subsection{Context and Motivation}

Generative modelling has undergone a paradigm shift in recent
years.  Early methods such as Generative Adversarial Networks (GANs)
relied on adversarial training and often suffered from mode collapse
and training instability.  Recent advances have introduced
likelihood‑based approaches that cast the generative task as the
reversal of a well‑defined forward process.  Among these, \emph{denoising
diffusion probabilistic models} have gained popularity for their
stability, expressivity and theoretical grounding.

Financial time series pose a particular challenge due to their
non‑stationarity and heavy tails.  Realistic simulation of stock
returns must reproduce stylised facts such as volatility clustering,
leptokurtic distributions and negligible autocorrelation at non‑zero
lags.  Traditional parametric models, including GARCH and
stochastic‑volatility formulations, capture some of these features but
often fail to reproduce the full distributional complexity observed in
practice.  Deep generative models promise to overcome these
limitations by directly learning the data manifold.

\subsection{Objectives of This Study}

This work has several intertwined objectives:

\begin{enumerate}[label=\arabic*., leftmargin=2em]
  \item To derive and prove key theoretical results in diffusion
    models, including the closed‑form expression for the marginal
    forward distribution and the rationale behind noise prediction.
  \item To explain and implement practical diffusion‑model variants
    such as DDPM, DDIM and classifier‑free guidance (CFG), assessing
    their relative strengths.
  \item To introduce Flow Matching as a continuous‑time alternative
    that learns velocity fields rather than stochastic scores, and to
    compare the two paradigms both mathematically and empirically.
  \item To apply these models to high‑frequency financial data, in
    particular the SPY ETF log‑returns, demonstrating how generative
    models can be used for synthetic data generation in finance.
  \item To provide a reproducible, modular implementation and a
    comprehensive evaluation using metrics such as SWD, autocorrelation
    functions, statistical moments and kurtosis.
\end{enumerate}

\subsection{Structure of the Report}

The remainder of this report is organised into seven additional
chapters, each in its own \TeX{} file.  Part~\ref{sec:introduction}
provides the context and sets up the problem.  Part~2 revisits the
theoretical foundations of diffusion models, deriving the closed‑form
marginal distribution and discussing noise prediction, variational
objectives, DDIM sampling and classifier‑free guidance.  Part~3
delves into the practical implementation of denoising diffusion
models, detailing the U‑Net architecture, training loops, sampling
algorithms and quantitative results.  Part~4 introduces Stable
Diffusion and DreamBooth, covering latent diffusion in the VAE
representation and personalised fine‑tuning via LoRA.  Parts 5 and 6
develop the theoretical and practical aspects of Flow Matching,
respectively, highlighting the deterministic ODE approach.  Part~7
presents an integrated evaluation and draws conclusions.  Part~8
contains appendices with detailed derivations, code listings,
additional figures and a mapping of the report contents to the
assignment specification.

\section{Methodology}
\label{sec:methodology}

\subsection{Data Preparation}

Our experiments utilise daily closing prices of the SPY ETF from
January~2010 to December~2023.  We compute log‑returns
$r_t = \log(p_t / p_{t-1})$ and standardise them globally to zero
mean and unit variance using the training set statistics.  The time
series is segmented into overlapping sequences of length $L=64$ for
the diffusion models and $L=100$ for Flow Matching.  We use a
sliding window with stride 1 to maximise the number of sequences.  A
chronological split divides the data into 90\% training and 10\%
testing.

\subsection{Model Architectures}

\paragraph{Diffusion Models.}  The core of the diffusion models is a
\emph{U‑Net} architecture tailored for noise prediction.  The U‑Net
comprises an encoder that progressively downsamples the image (or in
our case the one‑dimensional time series reshaped as a small image),
a bottleneck with self‑attention, and a decoder that upsamples back to
the original resolution.  Skip connections preserve high‑resolution
details.  Time embeddings are injected at each stage using
sinusoidal positional encodings.

\paragraph{Flow Matching.}  For Flow Matching we employ a
transformer‑based velocity field network.  Input sequences of
dimensionality $L \times F$ are projected into a hidden space and
processed by several multi‑head self‑attention layers.  Continuous
time information is encoded via sinusoidal embeddings and added to
every token.  The output projection yields a velocity vector for
each timestep.

\subsection{Training and Optimisation}

Diffusion models are trained by minimising a simplified MSE loss
between the predicted noise and the true noise at randomly sampled
timesteps.  We employ linear beta schedules with $T=1000$ steps and
optimise using AdamW with a learning rate of $10^{-4}$.  We train
for 50 epochs with batch size 128.  For DDIM sampling we use 50
inference steps.

Flow Matching training minimises the mean squared error between the
predicted velocity and the target velocity $x_1 - x_0$ at random
times $t \in [0,1]$.  We use batches of size 32 and train for 100
epochs with Adam optimiser (learning rate $10^{-4}$).  Euler
integration with 100 time steps approximates the ODE during
sampling.

\subsection{Evaluation Metrics}

We employ a range of metrics to assess the quality of the generated
time series.  The Sliced Wasserstein Distance (SWD) measures the
distance between high‑dimensional distributions by projecting
onto random directions and computing one‑dimensional Wasserstein
distances.  The autocorrelation function (ACF) quantifies temporal
dependencies; we compute both mean squared error (MSE) and mean
absolute error (MAE) between the real and generated ACFs.  Basic
statistical measures include mean, variance, skewness, kurtosis and
volatility.  Histograms and kernel density estimations (KDE) visualise
the distribution of returns.  A Q‑Q plot assesses tail behaviour.

\subsection{Scope and Limitations}

This study focuses on daily log‑returns of a single ETF.  While the
methods generalise to other assets and frequencies, caution should be
exercised when extrapolating to intraday data or multiple assets
without further model adaptations.  Our linear Flow Matching paths
provide a simple baseline; more sophisticated transport couplings may
yield improved results.  Finally, heavy‑tail reproduction remains
challenging, as we discuss in later sections.

\subsection{Notation}

Throughout the report we denote by $x_0$ the original clean sample,
$x_t$ the noised sample at diffusion step $t$, and $\epsilon$ the
standard Gaussian noise.  For Flow Matching we use $x_0$ for the
initial noise sample, $x_1$ for the target data sample and $x(t)$ for
the sample at continuous time $t \in [0,1]$.  Bold symbols denote
vectors and capital letters denote tensors or matrices.

\subsection{Reader Guidance}

Each part of the report can be read independently, but the
recommended order follows the enumeration above.  Mathematics
enthusiasts may appreciate diving into Part~2 early, while engineers
interested in implementation details may skip ahead to Parts~3–6.
Readers unfamiliar with machine learning or stochastic processes may
consult the appendices (Part~8) for a gentle primer on relevant
concepts.

\subsection{Historical Context}

The field of deep generative modelling has evolved rapidly over the
last decade.  Initial breakthroughs such as Variational Autoencoders
(VAEs) and Generative Adversarial Networks (GANs) demonstrated that
neural networks could learn to generate plausible images and other
structured data from latent codes.  VAEs maximise a variational
objective similar to that used in diffusion models but often produce
blurry samples due to the Gaussian prior imposed on the latent
space.  GANs achieve sharp samples via adversarial training but
suffer from instability and mode collapse.  Diffusion models marry
likelihood-based training with adversarial-like perceptual quality,
offering a robust alternative.  Flow-based models, normalising
flows and flow matching extend these ideas to invertible or
continuous-time dynamics.  In this report we build upon this
historical progression to compare diffusion and flow matching both
mathematically and empirically.

\subsection{Further Reading}

Readers interested in a deeper exploration of generative models may
consult the following references: Kingma and Welling’s original paper
on VAEs; Goodfellow et al.’s seminal work on GANs; Dinh et al.’s
RealNVP on normalising flows; Song and Ermon’s work on score-based
generative modelling; and recent surveys on diffusion models.  These
papers provide broader context for the techniques employed here and
highlight the vibrant research landscape in generative modelling.
